% Generated draft: requires mathematical and source review.
\documentclass[11pt]{article}
\usepackage[margin=25mm]{geometry}
\usepackage{fontspec}
\setmainfont{Latin Modern Roman}
\usepackage{amsmath,amssymb,mathtools,unicode-math}
\setmathfont{Latin Modern Math}
\usepackage{enumitem,needspace,xurl,hyperref}
\setlength{\parindent}{0pt}
\setlength{\parskip}{5pt}
\setlength{\emergencystretch}{3em}
\newcommand{\N}{\mathbb{N}}
\newcommand{\Z}{\mathbb{Z}}
\newcommand{\Q}{\mathbb{Q}}
\newcommand{\R}{\mathbb{R}}
\newcommand{\C}{\mathbb{C}}
\newcommand{\F}{\mathbb{F}}
\newcommand{\A}{\mathbb{A}}
\newcommand{\PP}{\mathbb{P}}
\newcommand{\E}{\mathbb{E}}
\newcommand{\eps}{\varepsilon}
\newcommand{\dd}{\,\mathrm{d}}
\begin{document}
\section*{UnsolvedMath ID: 700670}
\textbf{Source problem code:} catalogue entry
\subsection{Definitions and mathematical statement}
Fix a prime number \(p\). An abelian \(p\)-group is an abelian group \(G\) such that every element is annihilated by some power of \(p\). No cardinality bound on \(G\) is specified. Put \(R=\operatorname{End}_{\mathbb Z}(G)\), with multiplication given by composition, and regard \(G\) as a left \(R\)-module by evaluation: \(r\cdot g=r(g)\).

For a unital ring \(R\), a left \(R\)-module \(M\) is called slender if every \(R\)-linear map
\[
 f:\prod_{n=1}^{\infty}R\longrightarrow M
\]
satisfies \(f(e_n)=0\) for all but finitely many \(n\), where \(e_n\) has \(1_R\) in coordinate \(n\) and zero elsewhere. Thus the quantifiers are
\[
 \forall f\in\operatorname{Hom}_R\!\left(\prod_{n=1}^{\infty}R,M\right)\quad
 \exists F\subseteq\{1,2,\ldots\}\text{ finite}\quad
 \forall n\notin F,\quad f(e_n)=0.
\]

\textbf{Problem.} Characterize all abelian \(p\)-groups \(G\) that are slender as modules over \(\operatorname{End}_{\mathbb Z}(G)\) under the evaluation action. The characterization is sought for each prime \(p\).

The supplied source states that the standard group
\[
 B=\bigoplus_{n=1}^{\infty}\mathbb Z/p^n\mathbb Z
\]
is slender over its own endomorphism ring, attributing this fact to A. Mader. This is retained as supplied background, without independent verification. The source does not explicitly define module slenderness; the definition above states the convention used in this draft. The record also contains differing attribution signals: its statement begins with A. Mader, whereas its proposer metadata names S. V. Rychkov.
\subsection{Short English statement}
For each prime \(p\), characterize the abelian \(p\)-groups that are slender under the natural action of their endomorphism rings, using the explicit definition above. The source gives \(\bigoplus_{n\geq1}\mathbb Z/p^n\mathbb Z\) as an example. The supplied registry labels the problem open; that label has not been independently verified.
\subsection{Sources}
\begin{itemize}
\item Local UnsolvedMath registry, supplied complete record: permanent ID \(700670\), exact source problem code \texttt{catalogue entry}, title \emph{Kourovka Notebook Problem 10.55}; upstream revision \texttt{d371ee917ac9647fbb06cd6fc529deec6f3dda82}. The supplied status is \texttt{open}.
\item \emph{The Kourovka Notebook (Unsolved Problems in Group Theory)}, Problem 10.55, originating in the 10th issue (1986); the supplied background identifies the 21st issue edition (2026), arXiv:1401.0300v48, \url{https://arxiv.org/abs/1401.0300v48}. The supplied collection description names E. I. Khukhro and V. D. Mazurov as editors. These bibliographic details are supplied source data, not independently checked here.
\item A. Mader, contribution in \emph{Abelian Groups and Modules, Proc., Udine, 1984}, Springer, 1984, pp.~315--327, cited in the supplied statement for the assertion about the standard group \(B\). No contribution title is supplied.
\item Attribution qualification: the statement begins ``(A. Mader)'', while the registry's proposer field and background name S. V. Rychkov, with proposed year 1986. Both attributions are retained pending review.
\end{itemize}
\end{document}
